
\documentclass{article}
\usepackage{colortbl}
\usepackage{makecell}
\usepackage{multirow}
\usepackage{supertabular}

\begin{document}

\newcounter{utterance}

\centering \large Interaction Transcript for game `bbh', experiment `formal\_fallacies', episode 160 with Qwen3.5{-}4B.
\vspace{24pt}

{ \footnotesize  \setcounter{utterance}{1}
\setlength{\tabcolsep}{0pt}
\begin{supertabular}{c@{$\;$}|p{.15\linewidth}@{}p{.15\linewidth}p{.15\linewidth}p{.15\linewidth}p{.15\linewidth}p{.15\linewidth}}
    \# & \multicolumn{2}{c}{Player} && \multicolumn{2}{c}{Game Master} \\
    \hline

    \theutterance \stepcounter{utterance}  
    & & & \multicolumn{4}{p{0.6\linewidth}}{
        \cellcolor[rgb]{0.9,0.9,0.9}{
            \makecell[{{p{\linewidth}}}]{
                \texttt{\tiny{[P1$\langle$GM]}}
                \texttt{Distinguish deductively valid arguments from formal fallacies.} \\
\\ 
\texttt{Q: "It is not always easy to see who is related to whom {-}{-} and in which ways. The following argument pertains to this question: To begin with, Lesley is a close friend of Fernando. Moreover, being a close friend of Fernando or a schoolmate of Lowell is sufficient for being a great{-}grandfather of Leroy. It follows that Lesley is a great{-}grandfather of Leroy."} \\
\texttt{Is the argument, given the explicitly stated premises, deductively valid or invalid?} \\
\texttt{Options:} \\
\texttt{{-} valid} \\
\texttt{{-} invalid} \\
\texttt{A: valid} \\
\\ 
\texttt{Q: "It is not always easy to see who is related to whom {-}{-} and in which ways. The following argument pertains to this question: Whoever is not a great{-}grandfather of Clyde is a stepbrother of Brian. Being an ancestor of Dana is sufficient for not being a great{-}grandfather of Clyde. We may conclude: Everyone who is an ancestor of Dana is a stepbrother of Brian, too."} \\
\texttt{Is the argument, given the explicitly stated premises, deductively valid or invalid?} \\
\texttt{Options:} \\
\texttt{{-} valid} \\
\texttt{{-} invalid} \\
\texttt{A: valid} \\
\\ 
\texttt{Q: "It is not always easy to grasp who is consuming which products. The following argument pertains to this question: Every infrequent user of Paul Mitchell shampoo is either a rare consumer of Nioxin shampoo or a loyal buyer of Caress soap, or both. No regular consumer of Lush soap is a rare consumer of Nioxin shampoo and, in the same time, a loyal buyer of Caress soap. It follows that whoever is an infrequent user of Paul Mitchell shampoo is not a regular consumer of Lush soap."} \\
\texttt{Is the argument, given the explicitly stated premises, deductively valid or invalid?} \\
\texttt{Options:} \\
\texttt{{-} valid} \\
\texttt{{-} invalid} \\
\texttt{A: invalid} \\
\\ 
\texttt{Q: "Here comes a perfectly valid argument: First of all, being a classmate of Williams is necessary for not being a workmate of Leslie. Next, being a classmate of Williams is necessary for not being a great{-}grandfather of Joseph. Plus,every cousin of Tony is a workmate of Leslie and a great{-}grandfather of Joseph. Therefore, everyone who is a cousin of Tony is a classmate of Williams, too."} \\
\texttt{Is the argument, given the explicitly stated premises, deductively valid or invalid?} \\
\texttt{Options:} \\
\texttt{{-} valid } \\
\texttt{{-} invalid} \\
\texttt{A:} \\
            }
        }
    }
     \\ \\

    \theutterance \stepcounter{utterance}  
    & \multicolumn{4}{p{0.6\linewidth}}{
        \cellcolor[rgb]{0.561,0.737,0.561}{
            \makecell[{{p{\linewidth}}}]{
                \texttt{\tiny{[P1$\rangle$GM]}}
                \texttt{To determine the validity of this argument, let's break down the premises and the conclusion using logical notation.} \\
\\ 
\texttt{**1. Symbolize the Premises:**} \\
\texttt{*   **Premise 1:** "Being a classmate of Williams is necessary for not being a workmate of Leslie."} \\
\texttt{    *   Let \$C\$ = Classmate of Williams.} \\
\texttt{    *   Let \$W\$ = Workmate of Leslie.} \\
\texttt{    *   "Necessary for" means if the result happens, the condition must have happened (\$R \textbackslash{}to C\$).} \\
\texttt{    *   Structure: \$\textbackslash{}neg W \textbackslash{}to C\$ (If not a workmate, then a classmate).} \\
\texttt{*   **Premise 2:** "Being a classmate of Williams is necessary for not being a great{-}grandfather of Joseph."} \\
\texttt{    *   Let \$G\$ = Great{-}grandfather of Joseph.} \\
\texttt{    *   Structure: \$\textbackslash{}neg G \textbackslash{}to C\$ (If not a great{-}grandfather, then a classmate).} \\
\texttt{*   **Premise 3:** "Every cousin of Tony is a workmate of Leslie and a great{-}grandfather of Joseph."} \\
\texttt{    *   Let \$T\$ = Cousin of Tony.} \\
\texttt{    *   Structure: \$T \textbackslash{}to (W \textbackslash{}land G)\$ (If a cousin, then a workmate AND a great{-}grandfather).} \\
\\ 
\texttt{**2. Analyze the Conclusion:**} \\
\texttt{*   **Conclusion:** "Everyone who is a cousin of Tony is a classmate of Williams, too."} \\
\texttt{    *   Structure: \$T \textbackslash{}to C\$.} \\
\\ 
\texttt{**3. Evaluate the Logic:**} \\
\texttt{*   From **Premise 3**, we know that anyone who is a cousin of Tony (\$T\$) is a **workmate** (\$W\$) and a **great{-}grandfather** (\$G\$).} \\
\texttt{*   Therefore, for a cousin of Tony, the condition "not being a workmate" (\$\textbackslash{}neg W\$) is **false**, and the condition "not being a great{-}grandfather" (\$\textbackslash{}neg G\$) is **false**.} \\
\texttt{*   Now look at **Premise 1** (\$\textbackslash{}neg W \textbackslash{}to C\$). Since the antecedent (\$\textbackslash{}neg W\$) is false for our subject, the implication tells us nothing about the consequent (\$C\$). Just because the condition for being a classmate wasn't triggered by *not* being a workmate, it doesn't mean they *are* a classmate. They could be a classmate for other reasons, or they could be neither.} \\
\texttt{*   Similarly, **Premise 2** (\$\textbackslash{}neg G \textbackslash{}to C\$) has a false antecedent (\$\textbackslash{}neg G\$), so it provides no information about whether they are a classmate.} \\
\texttt{*   The premises establish that cousins are \$W\$ and \$G\$. The rules only tell us what happens if someone is *not*} \\
            }
        }
    }
    & & \\ \\

    \theutterance \stepcounter{utterance}  
    & & & \multicolumn{2}{p{0.3\linewidth}}{
        \cellcolor[rgb]{0.9,0.9,0.9}{
            \makecell[{{p{\linewidth}}}]{
                \texttt{\tiny{[GM$|$GM]}}
                \texttt{Target: invalid} \\
            }
        }
    }
    & & \\ \\

    \theutterance \stepcounter{utterance}  
    & & & \multicolumn{2}{p{0.3\linewidth}}{
        \cellcolor[rgb]{0.9,0.9,0.9}{
            \makecell[{{p{\linewidth}}}]{
                \texttt{\tiny{[GM$|$GM]}}
                \texttt{game\_result = LOSE} \\
            }
        }
    }
    & & \\ \\

\end{supertabular}
}

\end{document}
